\chapter{总结与展望}

\section{研究总结}

面向电动汽车全生命周期高可靠运行的需求，本文针对实际复杂工况下观测数据统计特性演化产生的深度网络推断偏差和泛化失效的问题，建立了基于无偏不确定性量化的~IPHM~研究主线。
% 
本文围绕动力电池可信状态估计、动力电池可泛化退化建模和旋转机械抗幻觉故障诊断三个问题展开研究，提出三种数据驱动的深度~PHM~方法，并基于主流开源锂离子电池循环老化数据集和旋转机械故障数据集开展广泛验证，所提出方法的整体性能优于多种对比方法。
主要结论如下：

（1）提出基于层次化随机时空~Transformer~的锂离子电池健康状态可信估计方法。

（2）提出基于拓扑约束下不确定性驱动特征调弈的锂离子电池分布外鲁棒退化建模方法。

（3）提出基于扩散增强可信多模态大语言模型的旋转机械抗幻觉故障诊断方法。

\section{工作展望}

本文在锂离子电池可信健康状态估计、锂离子电池分布外退化模式鲁棒退化建模和旋转机械抗幻觉故障诊断等方面开展广泛而深入的研究，取得了一系列成果。
所提出的方法在可解释性、与研究对象物理机理的结合性、规模性和可拓展性、任务多样性以及在应用场景中的自动化程度等方面仍有待进一步深入研究。
基于此，未来拟从以下几个方面展开深入研究：

（1）\textbf{探索计及不确定性的~PINN~及其在电池退化建模问题中的应用}。
% 
以~PINN~为代表的数据驱动-机理融合建模范式在电池退化建模问题中取得了良好的效果，对构建可解释电池管理体系有重要意义。PINN~构建过程中不可避免地受到物理模型简化引起的认知不确定性影响，PINN~训练的过程同样受到数据测量噪声和工况随机性引起的固有不确定性影响。量化这种不确定性，并利用此量化结果进行控制方程有助于模型对齐退化观测量和拟合电池内部电化学反应过程。
本文的未来工作将深入研究针对这种新范式的不确定性量化、校准和应用方法，以为下游的监控、决策和设计环节提供更准确可靠的电池模型。

（2）\textbf{探索面向跨化学体系、跨相态、跨尺度、跨应用场景的大规模电池基础模型}。
% 
本文广泛地研究了不同充放电循环协议、不同工作温度、不同具体电极材料的应用在车辆动力系统中的锂离子电池的退化建模问题。近年来随着电池技术的快速发展，对诸如钠离子电池等采用不同电荷载体体系的新型电池和以固态电池为代表的不同形态的电池技术的研究越发充分；大规模电池成组成包甚至承担结构件功能的相关研究方兴未艾，电池在垂直起降飞行器、卫星、储能电站和不间断电源等基础设施中的应用得到了广泛的关注，依据循环数据提取跨化学体系、跨相态、跨尺度和跨应用场景的通用动力学特征极具挑战。
未来工作拟开发电池基础模型，依托已有异构数据开展自监督训练，进一步增强数据驱动~PHM~模型的跨域迁移能力，解决新化学体系、新相态、大尺度和新应用场景下少样本建模问题。

（3）以（2）中研究为基础，\textbf{探索电池基础模型在多个电池~PHM~任务中的应用能力}。
% 
本文针对锂离子电池主要以~SOH和~RUL~为指标讨论退化建模问题，锂离子电池的其他状态量，如~SOC~和~SOE~等，的估计问题同样是~PHM~的研究对象，这些状态估计任务在形式上和所研究的长期性能退化问题非常相近；但在输入形式、潜在反应过程和对特征的依赖性上与所研究的问题存在差异。
除此之外，电池模组或电池包内部各电芯间的均衡性评估指标估计、电池包的整体性能状态估计以及电池包级的故障诊断是应用中的重要关切，同样是~PHM~的关键任务。
开发针对特定任务的深度~PHM~模型面临可扩展性问题，具体表现为单个模型容量有限、输入信号难以共享和异构模型间难以集成等。
未来工作拟深入分析此挑战，研究如何设计统一的、基于电池基础模型实现多状态量多任务的统一管理。

（4）\textbf{探索面向工业自主闭环运维的故障诊断基础模型和代理系统}。
% 
本文提出的故障诊断基础模型已经具备具备从复杂时序信号识别异常、定位故障模式并生成维护建议的功能，但实施维护仍需要人工参与。工业设备数量和种类快速增长，高精密设备和高危环境下设备的维护需求与日俱增，依赖人工介入的运维模式存在响应滞后和维护成本高昂的挑战。
未来工作拟结合已有~PHM~基础模型和智能体技术，构建感知、诊断、决策和执行四位一体的自主闭环运维体系。具体地，构建包含信号采集预处理智能体、手册和日志检索智能体、故障诊断和决策智能体、运维动作编排智能体、基于机械臂或轮式机器人的智能实体和人机交互智能体等在内的多智能体协同框架，实现无需人工干预的故障自主排查、隔离和修复。
